\documentclass[11pt,a4paper]{article}
\usepackage[spanish,es-nodecimaldot]{babel}
\spanishdeactivate{<>}
\usepackage[utf8]{inputenc}
\usepackage[T1]{fontenc}
\usepackage{lmodern}
\usepackage{geometry}
\usepackage{xcolor}
\usepackage{hyperref}
\usepackage{booktabs}
\usepackage{tabularx}
\usepackage{array}
\usepackage{longtable}
\usepackage{enumitem}
\usepackage{tikz}
\usetikzlibrary{arrows.meta,positioning,shapes.geometric}
\usepackage{fancyhdr}
\usepackage{titlesec}
\usepackage{microtype}

\geometry{margin=2.1cm,headheight=14pt}
\definecolor{navy}{HTML}{17324D}
\definecolor{blue}{HTML}{246B9F}
\definecolor{teal}{HTML}{167A72}
\definecolor{orange}{HTML}{B46118}
\definecolor{lightblue}{HTML}{EAF3F8}
\definecolor{lightgray}{HTML}{F3F5F7}
\definecolor{darkgray}{HTML}{46515C}
\hypersetup{colorlinks=true,linkcolor=navy,urlcolor=blue,pdftitle={Guía de estudio: supuestos prácticos A1 informática}}
\pagestyle{fancy}\fancyhf{}
\lhead{\textcolor{darkgray}{Guía de estudio A1 informática}}
\rhead{\textcolor{darkgray}{Supuestos prácticos}}
\cfoot{\thepage}
\titleformat{\section}{\Large\bfseries\color{navy}}{\thesection.}{0.6em}{}
\titleformat{\subsection}{\large\bfseries\color{blue}}{\thesubsection}{0.6em}{}
\titleformat{\subsubsection}{\normalsize\bfseries\color{teal}}{\thesubsubsection}{0.6em}{}
\setlist[itemize]{leftmargin=1.25em,itemsep=2pt,topsep=3pt}
\setlist[enumerate]{leftmargin=1.5em,itemsep=2pt,topsep=3pt}
\newcommand{\key}[1]{\textbf{\textcolor{navy}{#1}}}
\newcommand{\tagtopic}[1]{\textcolor{white}{\fcolorbox{blue}{blue}{\textbf{#1}}}}
\newcommand{\boxnote}[1]{\begin{center}\fcolorbox{blue}{lightblue}{\parbox{0.92\textwidth}{#1}}\end{center}}
\newcolumntype{Y}{>{\raggedright\arraybackslash}X}

\begin{document}
\begin{titlepage}
\centering\vspace*{1.5cm}
{\Large\textcolor{darkgray}{PREPARACIÓN DEL CUERPO SUPERIOR DE SISTEMAS Y TI}}\\[1cm]
{\Huge\bfseries\textcolor{navy}{Guía transversal de estudio}}\\[0.25cm]
{\LARGE Supuestos prácticos A1 de informática}\\[1.4cm]
{\large Temario desarrollado + plan de estudio hasta el 16 de octubre de 2026}\\[2cm]
\begin{tikzpicture}[node distance=7mm]
\tikzset{b/.style={draw=blue,fill=lightblue,rounded corners,align=center,text width=3.1cm,minimum height=10mm}}
\node[b] (j) {Derecho y\\Administración};
\node[b,right=of j] (a) {Arquitectura y\\tecnología};
\node[b,right=of a] (s) {Seguridad y\\datos};
\node[b,below=of a] (p) {Planificación,\\contratación y examen};
\draw[-stealth,thick,blue] (j)--(a);\draw[-stealth,thick,blue] (a)--(s);\draw[-stealth,thick,blue] (j)--(p);\draw[-stealth,thick,blue] (s)--(p);
\end{tikzpicture}
\vfill
\begin{minipage}{0.84\textwidth}\small\centering
Documento elaborado a partir de los supuestos sobre ciclo presupuestario, Plataforma Digital Común de pensiones y Plataforma Digital de Cooperación entre Administraciones.\\[3pt]
\textbf{Fecha de referencia del plan: 24 de septiembre--16 de octubre de 2026.}
\end{minipage}
\end{titlepage}

\tableofcontents\newpage

\section{Cómo utilizar esta guía}

Los tres ejercicios no exigen memorizar tecnologías aisladas. Exigen demostrar que se sabe convertir un problema jurídico-organizativo en una solución de información \key{integral, segura, interoperable, implantable y justificable}. Los patrones comunes son:
\begin{itemize}
\item identificar actores, competencias, procesos y documentos;
\item separar dominios funcionales y responsabilidades;
\item diseñar arquitectura lógica y física;
\item asegurar identidad, autorización, trazabilidad, continuidad y protección de datos;
\item priorizar un MVP y justificar recursos, costes y calendario;
\item incluir datos, transparencia, accesibilidad e innovación sin comprometer la legalidad.
\end{itemize}

\subsection{Método de respuesta}
\begin{enumerate}
\item Comenzar con alcance, supuestos y objetivos.
\item Repetir la numeración exacta del enunciado.
\item En cada módulo indicar funciones, actores, datos, interfaces, controles y evidencias.
\item Dibujar al menos un diagrama lógico y uno físico o de proceso.
\item Justificar las decisiones frente a alternativas, no enumerar productos.
\item Terminar con riesgos, implantación, indicadores y conclusión ejecutiva.
\end{enumerate}

\boxnote{\key{Plantilla mental}: problema institucional $\rightarrow$ procesos $\rightarrow$ dominios $\rightarrow$ datos $\rightarrow$ integraciones $\rightarrow$ seguridad $\rightarrow$ arquitectura $\rightarrow$ MVP $\rightarrow$ evolución.}

\section{Temario I. Marco jurídico y Administración digital}

\subsection{1. Organización administrativa y cooperación}

La Ley 40/2015 estructura las relaciones interadministrativas sobre cooperación, colaboración y coordinación. Los órganos de cooperación pueden ser multilaterales o bilaterales y reunir AGE, comunidades autónomas, Ceuta, Melilla y, en determinados casos, entidades locales.

La Plataforma de Cooperación se justifica por el artículo 142.c) y la disposición adicional trigésima: debe facilitar reuniones, documentos, decisiones, acuerdos y trabajo electrónico de Conferencias Sectoriales y otros órganos. Cada Conferencia tiene reglamento interno, debe inscribirse en el registro correspondiente y puede contar con Comisión Sectorial y Grupos de trabajo.

\key{Qué estudiar}: artículos 142, 145--154 de la Ley 40/2015; funciones de Presidencia, Secretaría, Comisión Sectorial y Grupos de trabajo; convocatoria, orden del día, acuerdos, recomendaciones, planes conjuntos, actas y certificaciones.

\subsection{2. Procedimiento administrativo electrónico}

La Ley 39/2015 regula identificación, firma, registros, expediente, notificaciones, archivo y actuación administrativa automatizada. La Ley 40/2015 regula el funcionamiento del sector público, relaciones electrónicas, interoperabilidad y seguridad.

Un expediente electrónico debe conservar contexto, documentos, metadatos, versiones, firmas, sellos de tiempo, estados, responsables y evidencias. La plataforma debe diferenciar borrador, trámite, aprobación, publicación, archivo y expurgo.

\key{Respuesta de examen}: relacionar cada función con el acto administrativo que soporta. Una convocatoria no es solo un evento de calendario: es una comunicación formal con destinatarios, orden del día, documentos, fecha y evidencia de remisión.

\subsection{3. Transparencia, reutilización y accesibilidad}

La Ley 19/2013 exige publicidad activa y acceso a la información, con límites por protección de datos, seguridad, secreto comercial o deliberaciones reservadas. Los datos abiertos deben publicarse en formatos reutilizables, con metadatos, actualización y licencia.

La accesibilidad debe aplicarse desde el diseño: WCAG, lectura clara, teclado, contraste, lectores de pantalla, adaptación móvil, subtítulos y alternativas para personas con discapacidad o brecha digital.

\subsection{4. Interoperabilidad}

El Esquema Nacional de Interoperabilidad orienta estándares, expediente, documento, firma, archivo, reutilización y modelos de datos. Conviene citar directorios, intercambio de datos, APIs, formatos estructurados, metadatos, identificadores comunes y catálogo de servicios.

\subsection{5. Firma, identidad y archivo}

El Reglamento eIDAS y la normativa española de identificación y firma permiten distinguir autenticación, firma, sello electrónico, sello de tiempo y validación. No todo acceso necesita firma, pero sí los actos que expresan voluntad o certifican un resultado.

El archivo debe garantizar integridad, disponibilidad, autenticidad, conservación, migración y trazabilidad. Una grabación o transcripción puede ser evidencia auxiliar, pero no sustituye automáticamente el acta firmada.

\section{Temario II. Dominio presupuestario y financiero}

\subsection{6. Ciclo presupuestario}

El ciclo comprende planificación plurianual, objetivos de estabilidad y deuda, límite de gasto no financiero, directrices, propuestas de órganos gestores, consolidación, elaboración del proyecto, aprobación gubernamental, tramitación parlamentaria, enmiendas, Ley, distribución de créditos, ejecución, liquidación y rendición.

\subsection{7. Clasificación y datos presupuestarios}

El presupuesto se organiza por dimensiones orgánica, funcional o por programas y económica. El maestro presupuestario debe controlar ejercicios, centros, programas, aplicaciones, capítulos, conceptos, fuentes de financiación y relaciones entre estructuras.

\subsection{8. Ejecución del gasto}

Las fases son retención de crédito, autorización, compromiso o disposición, reconocimiento de la obligación, propuesta/ordenación del pago y pago. La aplicación debe impedir secuencias inválidas, importes superiores, imputaciones sin crédito y actuaciones sin expediente.

\subsection{9. Contratación pública}

La contratación se integra con el gasto: necesidad, memoria, pliegos, certificado de crédito, autorización, licitación, ofertas, valoración, adjudicación, firma, compromiso, ejecución, factura, recepción, obligación y pago. Hay que conocer procedimientos abierto, restringido, negociado justificado, menor, emergencia, acuerdos marco, sistemas dinámicos, contratos de obras, suministros, servicios y concesiones.

\subsection{10. Contabilidad, control y tesorería}

El órgano contable registra documentos de cada fase, realiza conciliaciones y participa en cierre, liquidación y Cuenta General. La fiscalización previa revisa legalidad, crédito y trámites; el control financiero y las auditorías revisan fiabilidad, cumplimiento, economía, eficacia y eficiencia. Tesorería ejecuta pagos y concilia con bancos.

\key{Idea común}: separación de funciones. El gestor propone, el fiscalizador controla, la autoridad autoriza, el contable registra y Tesorería paga.

\section{Temario III. Pensiones y Plataforma Digital Común}

\subsection{11. Modelo funcional de la PDC}

La Ley 12/2022 impulsa planes de pensiones de empleo, fondos de promoción pública abiertos y planes simplificados. El Real Decreto 885/2022 y el Real Decreto 668/2023 desarrollan el marco. La PDC da información común a partícipes, beneficiarios, promotores, gestoras, depositarias y órganos de supervisión.

La PDC inicial no gestiona transacciones individuales: gestoras y depositarias son las fuentes de datos mediante ficheros normalizados y EDITRAN. La plataforma recibe, valida, normaliza, conserva el histórico y ofrece consultas según perfil.

\subsection{12. Dominios PDC}
\begin{itemize}
\item ingesta EDITRAN, certificados, ACK, errores y cuarentena;
\item normalización, calidad, duplicados y conciliación;
\item fondos, planes, entidades, promotores y comisiones;
\item movimientos, aportaciones, derechos consolidados y económicos;
\item consultas ciudadanas y de órganos supervisores;
\item identidad, roles, representantes, apoderamientos y ámbitos;
\item informes, supervisión, reclamaciones, auditoría y trazabilidad.
\end{itemize}

\subsection{13. Modelo de autorización PDC}

RBAC determina funciones por rol; ABAC añade entidad, plan, fondo, relación jurídica, vigencia y tipo de información. Un partícipe solo consulta sus datos; una gestora consulta sus envíos; una comisión consulta su ámbito; DGOSS y DGSFP acceden según competencia. La autenticación puede apoyarse en Cl@ve y PGIS, certificados de entidad y mecanismos corporativos.

\section{Temario IV. Cooperación y Plataform@}

\subsection{14. Modelo de espacios}

Plataform@ debe representar una jerarquía de plataforma, órgano, Conferencia, Comisión, Grupo de trabajo y reunión. Cada espacio tiene miembros, administradores delegados, reglamento, plantillas, workflows, niveles de acceso, ciclo de vida documental y reglas de publicación.

\subsection{15. Reuniones y documentos}

El expediente de reunión cubre convocatoria, iniciativa, orden del día, documentación, asistentes, quórum, deliberación, votación, acuerdo o recomendación, acta, firma, certificación, publicación y seguimiento. El Secretario prepara, convoca, recibe comunicaciones, redacta actas, certifica y custodia.

\subsection{16. Votación y delegación}

Debe parametrizarse quórum, mayoría, voto nominal o reservado, peso de representación, voto particular y reglas del órgano. La delegación debe ser expresa, autenticada, limitada, revocable, trazable y compatible con la prohibición de doble voto. El sistema calcula y documenta; no sustituye la validación jurídica.

\subsection{17. Videoconferencia}

Debe asociarse a una reunión, controlar acceso y asistentes, permitir moderación, registrar incidencias y separar grabación/transcripción del acta. La transcripción automática es un borrador revisable. La emisión pública requiere autorización y revisión de confidencialidad.

\subsection{18. Conocimiento jurídico y cooperación}

El catálogo debe integrar convenios, legislación AGE/CCAA, reparto competencial y jurisprudencia constitucional. Debe tener fuente, vigencia, relaciones, materias, búsqueda, enlaces oficiales y control de actualización.

\section{Temario V. Arquitectura de sistemas y aplicaciones}

\subsection{19. Arquitectura lógica}

Separar presentación, identidad, dominios funcionales, integración, datos, documentos, analítica y operación. Las aplicaciones deben tener límites claros y comunicarse por APIs versionadas, eventos o procesos ETL. Identificar siempre qué sistema es fuente, qué sistema transforma y cuál publica.

\subsection{20. Arquitecturas de aplicación}

La opción generalmente defendible es un \key{monolito modular evolucionable} o una arquitectura de servicios de dominio. Microservicios solo cuando aporten escalabilidad, autonomía o aislamiento reales. Patrones útiles: máquina de estados, workflow, RBAC/ABAC, auditoría inmutable, idempotencia, versionado, correlación de eventos y separación operacional/analítica.

\subsection{21. Integración}

REST para consultas y comandos, mensajería para eventos, colas para cargas, ETL para históricos y APIs de terceros. Deben definirse contratos, versionado, timeouts, reintentos, idempotencia, circuit breaker, trazas y tratamiento de errores.

\subsection{22. Gestión documental}

Repositorio de objetos con metadatos, hash, versión, firma, sello de tiempo, clasificación, retención y control de descarga. La base relacional conserva metadatos y relaciones; el contenido pesado se almacena en objeto; la búsqueda usa índice separado.

\subsection{23. Interfaz y experiencia}

Diseño responsive, accesible y móvil; lenguaje claro; navegación por tareas; mensajes de error accionables; internacionalización; soporte de asistencia; compatibilidad multiplataforma y pruebas con usuarios reales.

\section{Temario VI. Infraestructura, nube y operación}

\subsection{24. Arquitectura física}

Esquema habitual: usuarios, CDN/WAF/balanceador, DMZ, frontend, API Gateway, servicios, bus, bases de datos, repositorio documental, analítica, monitorización y copias. Entornos separados: desarrollo, integración, pruebas, preproducción, producción y recuperación.

\subsection{25. Nube frente a CPD}

La nube aporta elasticidad, rapidez y servicios gestionados; el CPD aporta control, soberanía e integración con sistemas existentes. Una decisión madura suele ser híbrida. Justificar residencia, cifrado, claves, portabilidad, SLA, auditoría, reversibilidad, coste total y continuidad.

\subsection{26. Dimensionamiento}

Calcular usuarios nominales y concurrentes, crecimiento, almacenamiento transaccional y documental, multimedia, red, CPU, memoria, réplicas, backups, RPO/RTO y picos de campaña. Separar datos operativos de BI. No dimensionar solo con el volumen actual: incluir versiones, copias, logs y retención.

\subsection{27. DevSecOps y operación}

CI/CD, Git, ramas protegidas, revisión de código, pruebas automatizadas, SAST/DAST, dependencias, artefactos, infraestructura como código, observabilidad, logs centralizados, métricas, trazas, SIEM, gestión de vulnerabilidades y rollback.

\section{Temario VII. Seguridad, ENS y protección de datos}

\subsection{28. Dimensiones y categorización ENS}

Analizar confidencialidad, integridad, disponibilidad, autenticidad y trazabilidad. Presupuesto y cooperación suelen exigir integridad/trazabilidad muy altas; pensiones exige confidencialidad e integridad altas; disponibilidad se valora según servicio y plazos. La categoría debe derivarse de análisis de riesgos y responsabilidades, no afirmarse sin justificar.

\subsection{29. RGPD y privacidad}

Principios: licitud, finalidad, minimización, exactitud, limitación de conservación, integridad/confidencialidad y responsabilidad proactiva. Aplicar privacidad desde el diseño, DPIA cuando proceda, registro de tratamientos, derechos, brechas, encargados, transferencias y anonimización para datos estadísticos.

\subsection{30. Control de acceso}

Identidad federada, MFA en perfiles privilegiados y actos críticos, mínimo privilegio, segregación de funciones, altas/bajas, revisión periódica, acceso temporal y trazabilidad. Separar identidad, rol y ámbito: una persona puede tener roles distintos en órganos diferentes.

\subsection{31. Medidas técnicas y organizativas}
\begin{itemize}
\item cifrado, gestión de claves, firma, sellado y hash;
\item segmentación, WAF, bastionado, EDR, parcheado y pruebas de penetración;
\item logs inmutables, SIEM, alertas y respuesta a incidentes;
\item copias inmutables, alta disponibilidad y recuperación;
\item auditorías, formación, política de seguridad y análisis de riesgos;
\item continuidad, simulacros, RTO/RPO y operación manual de contingencia.
\end{itemize}

\section{Temario VIII. Datos, transparencia e innovación}

\subsection{32. Gobierno del dato}

Definir propietarios, responsables, catálogo, glosario, calidad, linaje, clasificación, ciclo de vida y controles de acceso. Separar zona raw, curada y de consumo en un lake o lakehouse. Documentar origen, periodicidad, transformación y finalidad.

\subsection{33. Analítica y cuadros de mando}

Diseñar indicadores antes de elegir herramienta. Presupuesto: crédito, compromiso, obligación y pago. PDC: cargas, errores, derechos y cobertura. Plataform@: reuniones, asistencia, acuerdos y cumplimiento. Indicadores deben tener definición, fórmula, fuente, periodicidad, propietario y límites de interpretación.

\subsection{34. Transparencia y datos abiertos}

Publicar solo lo autorizado, en formatos reutilizables, con metadatos, actualización, licencia y control de versiones. Anonimizar o agregar cuando exista riesgo de identificación. El portal público debe consumir una vista de publicación, nunca exponer directamente la base transaccional.

\subsection{35. Inteligencia artificial responsable}

Casos defendibles: clasificación documental, extracción de metadatos, detección de anomalías, predicción de carga, búsqueda semántica y transcripción asistida. Siempre con supervisión humana, explicabilidad, evaluación de sesgos, control de datos, registro de versiones y posibilidad de no automatizar una decisión jurídicamente relevante.

\section{Temario IX. Planificación, recursos y contratación}

\subsection{36. Cronograma y MVP}

Fases: inicio y gobierno; requisitos y arquitectura; diseño; desarrollo iterativo; integraciones; pruebas; piloto; migración/formación; producción; soporte y evolución. El MVP debe cubrir el camino de mayor valor y riesgo: identidad, datos maestros, expediente, documentos, controles, auditoría y un flujo completo.

\subsection{37. Recursos}

Gobierno: comité, dirección y PMO. Funcional: propietarios de procesos, jurídico, datos, seguridad, archivo, transparencia y usuarios. Técnico: arquitectura, análisis, UX, desarrollo, integración, datos, QA, DevOps y operación. No olvidar soporte, formación, gestión del cambio y continuidad.

\subsection{38. Contratación pública TIC}

Justificar necesidad, insuficiencia de medios, objeto, lotes, solvencia, criterios de adjudicación, niveles de servicio, seguridad, propiedad intelectual, reutilización, transferencia y reversibilidad. El procedimiento abierto suele ser la opción ordinaria para servicios de desarrollo; valorar acuerdos marco o contratación centralizada si existen.

\subsection{39. Estimación económica}

Separar desarrollo, licencias, infraestructura, integración, seguridad, pruebas, formación, migración y explotación recurrente. Incluir contingencia y no ocultar costes de nube, almacenamiento, videoconferencia, soporte y evolución.

\section{Técnica para redactar en el examen}

\subsection{40. Esqueleto de respuesta}
\begin{enumerate}
\item Introducción: objetivo, alcance, supuestos y principios.
\item Modelo lógico: actores, dominios, datos, procesos y diagrama.
\item Tecnología: arquitectura, alternativas, infraestructura e integración.
\item Seguridad: impacto, ENS, RGPD, identidad, acceso, continuidad y auditoría.
\item Planificación: MVP, fases, hitos, equipos, metodología y costes.
\item Innovación/datos: valor, impacto, riesgos y evolución.
\item Conclusión: beneficios, riesgos residuales y decisión recomendada.
\end{enumerate}

\subsection{41. Frases de alto rendimiento}
\begin{itemize}
\item ``La solución separa la operación transaccional de la explotación analítica para proteger rendimiento e integridad.''
\item ``La autorización combina rol y ámbito; conocer la identidad no concede acceso universal.''
\item ``La parametrización permite adaptar reglamentos sin modificar código.''
\item ``La evidencia de cada actuación forma parte del expediente y se conserva con integridad y sello de tiempo.''
\item ``El MVP prioriza el flujo completo de mayor riesgo y deja la innovación avanzada para una evolución controlada.''
\end{itemize}

\subsection{42. Errores que restan puntos}
\begin{itemize}
\item Enumerar marcas sin justificar requisitos.
\item Proponer microservicios por moda o una única aplicación monolítica para todo.
\item Olvidar actores, responsabilidades, documentos o separación de funciones.
\item Confundir autenticación con autorización.
\item Publicar datos sin hablar de privacidad, anonimización y ciclo de vida.
\item Omitir operación, continuidad, soporte, formación y reversibilidad.
\item No responder exactamente a un apartado o no hacer visible el MVP.
\end{itemize}

\section{Calendario de estudio: 24 de septiembre--16 de octubre}

El plan supone aproximadamente 4--5 horas diarias: dos bloques de estudio, un bloque de recuperación activa y un ejercicio escrito. Si se dispone de menos tiempo, conservar siempre el repaso activo y los simulacros, reduciendo la lectura pasiva.

\subsection{Semana 1: base jurídica y modelo de respuesta}
\begin{center}
\begin{tabularx}{0.98\textwidth}{l l Y Y}\toprule
\textbf{Fecha} & \textbf{Tema} & \textbf{Trabajo principal} & \textbf{Salida obligatoria}\\\midrule
24/09 & Diagnóstico & Leer los tres enunciados; extraer actores, puntos y fechas. & Mapa común de 1 página.\\
25/09 & Ley 39/2015 & Expediente, identificación, firma, registro, notificación y archivo. & Esquema + respuesta de 30 min.\\
26/09 & Ley 40/2015 & Cooperación, órganos, Conferencias, Secretarías y decisiones. & Tabla de artículos y diagrama Plataform@.\\
27/09 & ENS, RGPD, ENI & Dimensiones, riesgos, medidas, privacidad e interoperabilidad. & Matriz de seguridad.\\
28/09 & Contratación TIC & Procedimientos, lotes, criterios, SLA, reversibilidad. & Relación valorada de 1,35 M€.\\
29/09 & Presupuesto I & Planificación, aprobación, clasificación y crédito. & Flujo presupuestario completo.\\
30/09 & Presupuesto II & Gasto, contratación, contabilidad, control y Tesorería. & Caso de expediente de gasto.\\\bottomrule
\end{tabularx}\end{center}

\subsection{Semana 2: dominios y arquitectura}
\begin{center}
\begin{tabularx}{0.98\textwidth}{l l Y Y}\toprule
\textbf{Fecha} & \textbf{Tema} & \textbf{Trabajo principal} & \textbf{Salida obligatoria}\\\midrule
01/10 & Arquitectura lógica & Dominios, límites, actores, datos y APIs. & Diagrama común reutilizable.\\
02/10 & Integración & REST, eventos, colas, ETL, errores, idempotencia y trazas. & Catálogo de interfaces.\\
03/10 & Arquitectura física & CPD, nube, híbrida, redes, entornos y continuidad. & Diagrama físico + dimensionamiento.\\
04/10 & Aplicaciones & Modularidad, workflows, documentos, búsqueda y UX. & Decisiones frente a alternativas.\\
05/10 & PDC pensiones & EDITRAN, PGIS, Cl@ve, cargas, histórico y perfiles. & Respuesta PDC 90 min.\\
06/10 & Plataform@ & Espacios, reuniones, actas, votación, videoconferencia y conocimiento. & Respuesta Plataform@ 90 min.\\
07/10 & Repaso activo & Recuperar de memoria los temas 1--23; corregir lagunas. & Fichas de errores.\\\bottomrule
\end{tabularx}\end{center}

\subsection{Semana 3: seguridad, datos y planificación}
\begin{center}
\begin{tabularx}{0.98\textwidth}{l l Y Y}\toprule
\textbf{Fecha} & \textbf{Tema} & \textbf{Trabajo principal} & \textbf{Salida obligatoria}\\\midrule
08/10 & Seguridad avanzada & ENS, RBAC/ABAC, MFA, auditoría, continuidad y DPIA. & Dos matrices de impacto.\\
09/10 & Datos y transparencia & Gobierno, lakehouse, BI, publicación, anonimización e IA. & Proyecto de datos por supuesto.\\
10/10 & MVP y cronograma & Priorización, dependencias, riesgos, pruebas y despliegue. & Tres MVP comparados.\\
11/10 & Recursos y contratación & Equipos, metodología, configuración y presupuesto. & Plan de proyecto completo.\\
12/10 & Simulacro 1 & Resolver un supuesto completo en tiempo limitado. & Documento corregido con rúbrica.\\
13/10 & Corrección 1 & Reescribir introducción, diagramas, seguridad y cierre. & Plantilla personal definitiva.\\
14/10 & Simulacro 2 & Resolver otro supuesto priorizando claridad y síntesis. & Segundo documento y lista de fallos.\\\bottomrule
\end{tabularx}\end{center}

\subsection{Recta final: 15 y 16 de octubre}
\begin{center}
\begin{tabularx}{0.98\textwidth}{l Y Y}\toprule
\textbf{Fecha} & \textbf{Plan} & \textbf{No hacer}\\\midrule
15/10 & Repaso de fichas, artículos, matrices, diagramas y frases de respuesta. Practicar introducción y conclusión. & No abrir temas nuevos ni leer documentación extensa.\\
16/10 & Repaso ligero: flujos, actores, ENS, MVP, arquitectura y contratación. Descanso y preparación material. & No realizar un simulacro completo ni estudiar de madrugada.\\\bottomrule
\end{tabularx}\end{center}

\section{Rúbrica personal de evaluación}

Después de cada simulacro, puntuar de 0 a 2 cada criterio:
\begin{center}
\begin{tabularx}{0.95\textwidth}{l c Y}\toprule
\textbf{Criterio} & \textbf{Puntuación} & \textbf{Pregunta de control}\\\midrule
Cobertura & 0--2 & ¿He respondido todos los apartados?\\
Modelo funcional & 0--2 & ¿Se entienden actores, procesos, datos y límites?\\
Arquitectura & 0--2 & ¿He justificado integración, tecnología y alternativas?\\
Seguridad & 0--2 & ¿He tratado ENS, privacidad, acceso, continuidad y auditoría?\\
Implantación & 0--2 & ¿El MVP, cronograma y recursos son realizables?\\
Datos e innovación & 0--2 & ¿Aporto valor sin olvidar legalidad y riesgos?\\
Expresión & 0--2 & ¿La respuesta es ordenada, legible y ejecutiva?\\
\bottomrule\end{tabularx}\end{center}

\boxnote{\key{Objetivo mínimo antes del 15 de octubre}: poder dibujar de memoria una arquitectura lógica, una física y un flujo de proceso; explicar ENS/RGPD; proponer un MVP; y redactar en menos de tres horas una respuesta completa, numerada y justificada.}

\section{Checklist final}
\begin{itemize}
\item ¿He identificado actores y responsabilidades?
\item ¿He distinguido información, operación, control y publicación?
\item ¿He definido datos maestros, documentos, histórico y auditoría?
\item ¿He explicado APIs, eventos, colas, errores e idempotencia?
\item ¿He justificado CPD, nube o híbrida con requisitos?
\item ¿He categorizado seguridad por dimensiones y no solo citado ENS?
\item ¿He separado autenticación, autorización y ámbito?
\item ¿He incluido protección de datos, accesibilidad y transparencia?
\item ¿He concretado MVP, cronograma, equipos y costes?
\item ¿He incluido operación, continuidad, soporte y reversibilidad?
\item ¿He reservado espacio a innovación responsable y sus riesgos?
\item ¿La conclusión responde al problema institucional inicial?
\end{itemize}

\section{Resumen de una página}

\begin{center}
\begin{tikzpicture}[font=\scriptsize,node distance=5mm]
\tikzset{b/.style={draw=blue,fill=lightblue,rounded corners,align=center,text width=4cm,minimum height=9mm},a/.style={-stealth,thick,blue}}
\node[b] (p) {1. Problema, alcance y actores};
\node[b,below=of p] (d) {2. Dominios, procesos, datos y documentos};
\node[b,below=of d] (i) {3. Integración, arquitectura lógica y física};
\node[b,below=of i] (s) {4. ENS, RGPD, identidad, acceso y trazabilidad};
\node[b,below=of s] (m) {5. MVP, cronograma, recursos y contratación};
\node[b,below=of m] (v) {6. Datos, transparencia, innovación y conclusión};
\draw[a] (p)--(d)--(i)--(s)--(m)--(v);
\end{tikzpicture}
\end{center}

La respuesta excelente mantiene coherencia entre el problema, el modelo funcional, la arquitectura, la seguridad y la implantación. No gana la propuesta con más tecnología, sino la que mejor demuestra criterio de Administración pública digital.

\end{document}
